\chapter{约束最优化方法}\label{chap:constrained-methods}
\section{非线性最小二乘问题}\label{sec:nls}

\subsection{问题描述}

\begin{definition}[非线性最小二乘问题]
给定训练集 $S = \{(x_i, y_i)\}_{i=1}^{m}$，需要估计近似函数 $f(x, \theta)$ 来模拟数据变化趋势。定义残量：
\[
r_i(\theta) = f(x_i, \theta) - y_i,\quad i = 1, 2, \ldots, m
\]

非线性最小二乘问题为：
\[
\min_{\theta \in \R^n}\ F(\theta) = \frac{1}{2}\sum_{i=1}^{m} r_i(\theta)^2 = \frac{1}{2}\|r(\theta)\|^2
\]
其中 $r: \R^n \to \R^m$ 是残量向量函数。若 $r$ 为线性函数，则退化为线性最小二乘。
\end{definition}

\subsection{梯度和Hessian矩阵}

设 $J(\theta)$ 是 $r(\theta)$ 的Jacobi矩阵：
\[
J(\theta) = \begin{pmatrix}
\frac{\partial r_1}{\partial \theta_1} & \cdots & \frac{\partial r_1}{\partial \theta_n} \\
\vdots & \ddots & \vdots \\
\frac{\partial r_m}{\partial \theta_1} & \cdots & \frac{\partial r_m}{\partial \theta_n}
\end{pmatrix} \in \R^{m \times n}
\]

\begin{property}[梯度和Hessian]
\begin{align*}
\nabla F(\theta) &= J(\theta)^T r(\theta) = \sum_{i=1}^{m} r_i(\theta) \nabla r_i(\theta) \\
\nabla^2 F(\theta) &= J(\theta)^T J(\theta) + S(\theta)
\end{align*}
其中 $S(\theta) = \sum_{i=1}^{m} r_i(\theta) \nabla^2 r_i(\theta)$ 包含残量的二阶信息。
\end{property}

\subsection{高斯-牛顿法（Gauss-Newton）}

\begin{conceptbox}[高斯-牛顿法原理]
在Hessian中忽略 $S(\theta)$ 项（假设残量较小或接近线性），得到近似Hessian $J^T J$。Newton迭代变为：
\[
\theta_{k+1} = \theta_k - (J_k^T J_k)^{-1} J_k^T r_k
\]
\end{conceptbox}

\begin{algobox}[高斯-牛顿法算法]
\textbf{Step 1}：解线性方程组 $J_k^T J_k \, d_k = -J_k^T r_k$ 得到方向 $d_k$

\textbf{Step 2}：更新 $\theta_{k+1} = \theta_k + d_k$（或加入线搜索）
\end{algobox}

\begin{notebox}[高斯-牛顿法的特点]
\begin{itemize}[itemsep=0.3em]
    \item 只需残量的一阶导数信息（Jacobi矩阵）
    \item $J^T J$ 至少半正定，保证下降方向
    \item 收敛速度取决于 $S(\theta)$ 的重要性：
    \begin{itemize}
        \item $S(\theta^*) = 0$：二阶收敛
        \item $S(\theta)$ 较小：Q-线性收敛
        \item $S(\theta)$ 较大：可能不收敛
    \end{itemize}
    \item 若 $J$ 不满秩，方法失效
    \item 带线搜索的版本称为\textbf{阻尼高斯-牛顿法}
\end{itemize}
\end{notebox}

\subsection{Levenberg-Marquardt方法}

\begin{conceptbox}[LM方法原理]
结合高斯-牛顿法和梯度下降法，引入信赖域思想：
\[
\min_d\ \frac{1}{2}\|J_k d + r_k\|^2 \quad \text{s.t.}\quad \|d\| \leq \Delta_k
\]

等价于求解：
\[
(J_k^T J_k + \lambda_k I) d = -J_k^T r_k
\]
其中 $\lambda_k > 0$ 为正则化参数。
\end{conceptbox}

\begin{notebox}[LM方法的特点]
\begin{itemize}[itemsep=0.3em]
    \item $\lambda_k$ 大时：接近梯度下降法（最速下降）
    \item $\lambda_k$ 小时：接近高斯-牛顿法
    \item $J^T J + \lambda I$ 正定，保证下降方向
    \item 是最常用的非线性最小二乘算法
    \item 通过调整 $\lambda$ 自适应地切换行为
\end{itemize}
\end{notebox}

\begin{example}[高斯-牛顿法迭代两次]\label{ex:gn}
用高斯-牛顿法拟合 $f(x, \theta) = \theta_1 e^{\theta_2 x}$，数据点 $(0, 2), (1, 3)$。初始 $\theta_0 = (1, 0)^T$。

\textbf{解}：残量 $r_1 = \theta_1 - 2$，$r_2 = \theta_1 e^{\theta_2} - 3$

\textbf{Jacobi矩阵}：
\[
J = \begin{pmatrix} 1 & 0 \\ e^{\theta_2} & \theta_1 e^{\theta_2} \end{pmatrix}
\]

\textbf{第0轮}：$\theta_0 = (1, 0)^T$
\[
r_0 = \begin{pmatrix} -1 \\ -2 \end{pmatrix},\quad
J_0 = \begin{pmatrix} 1 & 0 \\ 1 & 1 \end{pmatrix}
\]
\[
J_0^T J_0 = \begin{pmatrix} 2 & 1 \\ 1 & 1 \end{pmatrix},\quad
J_0^T r_0 = \begin{pmatrix} -3 \\ -2 \end{pmatrix}
\]
解 $J_0^T J_0 \, d_0 = -J_0^T r_0$ 得 $d_0 = (1, 1)^T$
\[
\theta_1 = \theta_0 + d_0 = (2, 1)^T
\]

\textbf{第1轮}：$\theta_1 = (2, 1)^T$
\[
r_1 = \begin{pmatrix} 0 \\ 2e - 3 \end{pmatrix} \approx \begin{pmatrix} 0 \\ 2.437 \end{pmatrix},\quad
J_1 = \begin{pmatrix} 1 & 0 \\ e & 2e \end{pmatrix} \approx \begin{pmatrix} 1 & 0 \\ 2.718 & 5.437 \end{pmatrix}
\]
解 $J_1^T J_1 \, d_1 = -J_1^T r_1$ 得新的方向 $d_1$，更新到 $\theta_2$。
\end{example}

\section{最优性条件回顾}\label{sec:optimality}

\subsection{无约束问题}

\begin{theorem}[无约束一阶必要条件]
设 $f$ 在 $x^*$ 处可微，若 $x^*$ 是局部最优解，则 $\nabla f(x^*) = 0$。
\end{theorem}

\begin{theorem}[无约束二阶必要条件]
设 $f$ 在 $x^*$ 处二阶可微，若 $x^*$ 是局部最优解，则 $\nabla f(x^*) = 0$ 且 $\nabla^2 f(x^*) \succeq 0$。
\end{theorem}

\begin{theorem}[无约束二阶充分条件]
设 $f$ 在 $x^*$ 处二阶可微，若 $\nabla f(x^*) = 0$ 且 $\nabla^2 f(x^*) \succ 0$，则 $x^*$ 是严格局部最优解。
\end{theorem}

\subsection{有约束问题——基本概念}

考虑约束优化问题：
\[
\min\ f(x)\quad \text{s.t.}\quad g_i(x) \leq 0,\ i \in \mathcal{I};\quad h_j(x) = 0,\ j \in \mathcal{E}
\]

\begin{definition}[可行域]
$\mathcal{X} = \{x \in \R^n \mid g_i(x) \leq 0,\ h_j(x) = 0,\ \forall i, j\}$
\end{definition}

\begin{definition}[起作用约束（Active Constraint）]
在可行点 $x$ 处，满足 $g_i(x) = 0$ 的不等式约束称为\textbf{起作用约束}。起作用约束集记为：
\[
\mathcal{A}(x) = \{i \in \mathcal{I} \mid g_i(x) = 0\}
\]
\end{definition}

\begin{definition}[可行方向]
方向 $d$ 在点 $x$ 处是\textbf{可行方向}，若存在 $\bar{\alpha} > 0$ 使得：
\[
x + \alpha d \in \mathcal{X},\quad \forall \alpha \in [0, \bar{\alpha}]
\]
\end{definition}

\begin{definition}[可行下降方向]
方向 $d$ 在点 $x$ 处是\textbf{可行下降方向}，若：
\begin{enumerate}[label=(\arabic*)]
    \item $d$ 是可行方向
    \item $\nabla f(x)^T d < 0$（下降方向）
\end{enumerate}
\end{definition}

\subsection{等式约束问题}

\begin{theorem}[等式约束的一阶条件]
设 $x^*$ 是等式约束问题的局部最优点，约束梯度线性无关（LICQ），则存在 $\lambda^*$ 使得：
\[
\nabla f(x^*) + \sum_{j \in \mathcal{E}} \lambda_j^* \nabla h_j(x^*) = 0
\]
即 $\nabla f(x^*)$ 是约束梯度的线性组合。
\end{theorem}

\begin{example}[验证最优性条件]
$\min\ f(x) = x_1 + x_2$，s.t. $x_1^2 + x_2^2 = 2$。验证 $x^* = (-1, -1)$ 是否满足一阶必要条件。

\textbf{解}：
\[
\nabla f = (1, 1)^T,\quad h(x) = x_1^2 + x_2^2 - 2 = 0,\quad \nabla h = (2x_1, 2x_2)^T
\]
在 $x^* = (-1, -1)$：
\[
\nabla h(x^*) = (-2, -2)^T
\]
需要找 $\lambda$ 使 $\nabla f + \lambda \nabla h = 0$：
\[
(1, 1)^T + \lambda(-2, -2)^T = 0 \Rightarrow 1 - 2\lambda = 0 \Rightarrow \lambda = \frac{1}{2}
\]
满足一阶必要条件。
\end{example}

\section{KKT条件回顾}\label{sec:ch7-kkt-review}

本章讨论约束优化的数值方法，这些方法通常需要判断当前迭代点是否满足最优性条件。
\begin{conceptbox}[KKT条件速查]
KKT条件的详细推导与证明见第~\ref{chap:constrained-theory}~章第~\ref{sec:ch6-kkt}~节。
\begin{itemize}
    \item \textbf{平稳性}：$\nabla f(x^*) + \sum \mu_i^* \nabla g_i(x^*) + \sum \lambda_j^* \nabla h_j(x^*) = 0$
    \item \textbf{原始可行性}：$g_i(x^*) \leq 0$，$h_j(x^*) = 0$
    \item \textbf{对偶可行性}：$\mu_i^* \geq 0$
    \item \textbf{互补松弛}：$\mu_i^* g_i(x^*) = 0$
\end{itemize}
\end{conceptbox}
\section{可行方向法}\label{sec:ch7-feasible-direction}

可行方向法的理论基础是最优性条件（见第~\ref{chap:constrained-theory}~章第~\ref{sec:ch6-feasible-direction}~节）：在局部最优点处不存在可行下降方向。因此，通过系统地寻找可行下降方向来迭代逼近最优解。

\subsection{基本思想}

可行方向法（Feasible Direction Method）在每次迭代中寻找一个\textbf{可行下降方向}，然后沿该方向进行一维搜索。

\begin{algobox}[可行方向法框架]
\textbf{迭代格式}：$x_{k+1} = x_k + \alpha_k d_k$

要求方向 $d_k$ 满足：
\begin{enumerate}[label=(\arabic*)]
    \item 可行性：$x_k + \alpha d_k \in \mathcal{X}$ 对充分小的 $\alpha > 0$
    \item 下降性：$\nabla f(x_k)^T d_k < 0$
\end{enumerate}
\end{algobox}

\subsection{线性约束情况}

考虑问题：
\[
\min\ f(x)\quad \text{s.t.}\quad Ax \leq b,\quad Ex = e
\]

\begin{theorem}[可行方向判定的充要条件]
设 $x$ 是可行解，$A_1 x = b_1$（起作用约束），$A_2 x < b_2$（不起作用约束）。则：
\begin{enumerate}[label=(\arabic*)]
    \item $d$ 是可行方向的充要条件：$A_1 d \leq 0$，$Ed = 0$
    \item 若还满足 $\nabla f(x)^T d < 0$，则 $d$ 是可行下降方向
\end{enumerate}
\end{theorem}

\textbf{求解可行下降方向}：转化为线性规划问题
\[
\min_{d, z}\ z\quad \text{s.t.}\quad \nabla f(x)^T d \leq z,\ A_1 d \leq 0,\ Ed = 0,\ -1 \leq d_j \leq 1
\]

\begin{notebox}[判定准则]
设最优解为 $(d^*, z^*)$：
\begin{itemize}[itemsep=0.3em]
    \item 若 $z^* < 0$：$d^*$ 是可行下降方向
    \item 若 $z^* = 0$：$x$ 是KKT点（找不到可行下降方向）
\end{itemize}
\end{notebox}

\subsection{非线性约束情况}

对于非线性约束 $g_i(x) \leq 0$，可行方向条件变为：
\[
\nabla g_i(x)^T d \leq 0,\quad \forall i \in \mathcal{A}(x)
\]

同样转化为线性规划求解 $d$ 和 $z$：
\[
\min_{d, z}\ z\quad \text{s.t.}\quad \nabla f(x)^T d \leq z,\ \nabla g_i(x)^T d \leq 0\ (i \in \mathcal{A}),\ -1 \leq d_j \leq 1
\]

\subsection{计算步骤}

\begin{algobox}[可行方向法计算步骤]
\textbf{Step 1}：给定初始可行点 $x_0$，允许误差 $\varepsilon > 0$，$k = 0$

\textbf{Step 2}：确定起作用约束集 $\mathcal{A}(x_k)$

\textbf{Step 3}：若 $\mathcal{A} = \emptyset$（内点）：
\begin{itemize}
    \item 若 $\|\nabla f(x_k)\| < \varepsilon$，停
    \item 否则取 $d_k = -\nabla f(x_k)$（最速下降），转Step 6
\end{itemize}

\textbf{Step 4}：若 $\mathcal{A} \neq \emptyset$（边界点），求解方向LP得 $(d_k^*, z^*)$

\textbf{Step 5}：若 $|z^*| < \varepsilon$，停（$x_k$ 近似KKT点）；否则令 $d_k = d_k^*$

\textbf{Step 6}：沿 $d_k$ 线搜索求 $\alpha_k$，更新 $x_{k+1} = x_k + \alpha_k d_k$

\textbf{Step 7}：$k \leftarrow k+1$，转Step 2
\end{algobox}

\begin{example}[可行方向法——迭代两次]\label{ex:fd}
$\min\ f(x) = 2x_1^2 + 2x_2^2 - 2x_1 x_2 - 4x_1 - 6x_2$
s.t. $x_1 + x_2 \leq 2$，$x_1 + 5x_2 \leq 5$，$x_1 \geq 0$，$x_2 \geq 0$。
取 $x^{(0)} = (0, 0)^T$，$\varepsilon = 0.001$。

\textbf{解}：$\nabla f = (4x_1 - 2x_2 - 4,\ 4x_2 - 2x_1 - 6)^T$。

记约束为 $g_1=x_1+x_2-2\leq0$，$g_2=x_1+5x_2-5\leq0$，$g_3=-x_1\leq0$，$g_4=-x_2\leq0$。

\textbf{第0轮}：$x^{(0)} = (0, 0)$。起作用约束为 $g_3,g_4$，故可行方向需满足 $d_1\geq0,d_2\geq0$。
\[
\nabla f(x^{(0)})=(-4,-6)^T
\]
方向LP为
\[
\min z\quad \text{s.t.}\quad -4d_1-6d_2\leq z,\ d_1\geq0,\ d_2\geq0,\ -1\leq d_j\leq1
\]
解得 $d^{(0)}=(1,1)^T$，$z=-10<0$。沿该方向线搜索：
\[
x^{(0)}+\alpha d^{(0)}=(\alpha,\alpha)^T,\quad
0\leq \alpha\leq \frac{5}{6}
\]
\[
f(\alpha,\alpha)=2\alpha^2-10\alpha
\]
在可行区间上取 $\alpha_0=5/6$，得
\[
x^{(1)}=\left(\frac{5}{6},\frac{5}{6}\right)^T
\]

\textbf{第1轮}：此时起作用约束为 $g_2$，且
\[
\nabla f(x^{(1)})=\left(-\frac{7}{3},-\frac{13}{3}\right)^T
\]
方向LP为
\[
\min z\quad \text{s.t.}\quad -\frac{7}{3}d_1-\frac{13}{3}d_2\leq z,\ d_1+5d_2\leq0,\ -1\leq d_j\leq1
\]
解得 $d^{(1)}=(1,-1/5)^T$，$z=-22/15<0$。沿该方向线搜索，得
\[
\alpha_1=\frac{55}{186},\quad
x^{(2)}=x^{(1)}+\alpha_1 d^{(1)}
=\left(\frac{35}{31},\frac{24}{31}\right)^T
\]

在 $x^{(2)}$ 处，只有 $g_2$ 起作用，且
\[
\nabla f(x^{(2)})=\left(-\frac{32}{31},-\frac{160}{31}\right)^T
=-\frac{32}{31}(1,5)^T
\]
取乘子 $\lambda_2=32/31\geq0$ 满足KKT条件。因此
\[
x^*=\left(\frac{35}{31},\frac{24}{31}\right)^T\approx(1.129,\ 0.774),\quad
f^*=-\frac{6882}{961}\approx -7.161
\]
\end{example}

\section{制约函数法（罚函数法与内点法）}\label{sec:ch7-penalty}

罚函数法与内点法的理论基础见第~\ref{chap:constrained-theory}~章第~\ref{sec:ch6-penalty}~节。本章侧重算法的具体实现与数值计算。

\subsection{外点罚函数法}

\begin{conceptbox}[外点罚函数法思想]
通过构造\textbf{惩罚函数} $P(x, \sigma)$，将约束问题转化为一系列无约束问题：
\[
\min_x\ P(x, \sigma) = f(x) + \sigma \cdot \text{penalty}(x)
\]
其中惩罚项在可行域内为零，在可行域外取很大的正值。
\end{conceptbox}

\textbf{常用惩罚函数}：

\textbf{等式约束}：$h_j(x) = 0$：
\[
P(x, \sigma) = f(x) + \frac{\sigma}{2}\sum_{j \in \mathcal{E}} h_j(x)^2
\]

\textbf{不等式约束}：$g_i(x) \leq 0$：
\[
g_i^+(x) = \max\{0, g_i(x)\},\quad
P(x, \sigma) = f(x) + \frac{\sigma}{2}\sum_{i \in \mathcal{I}} g_i^+(x)^2
\]

\textbf{一般形式}：
\[
P(x, \sigma) = f(x) + \frac{\sigma}{2}\left[\sum_{j \in \mathcal{E}} h_j(x)^2 + \sum_{i \in \mathcal{I}} (g_i^+(x))^2\right]
\]

\begin{algobox}[外点罚函数法算法]
\textbf{Step 1}：给定初始点 $x_0$，初始罚因子 $\sigma_0 > 0$，放大系数 $c > 1$，$\varepsilon > 0$，$k = 0$

\textbf{Step 2}：以 $x_k$ 为初值，求解 $\min_x P(x, \sigma_k)$，得 $x_{k+1}$

\textbf{Step 3}：若 $\sigma_k \cdot \text{penalty}(x_{k+1}) < \varepsilon$，停

\textbf{Step 4}：$\sigma_{k+1} = c \sigma_k$，$k \leftarrow k+1$，转Step 2
\end{algobox}

\begin{theorem}[外点罚函数法收敛性]
设 $\sigma_k \to +\infty$，$x_k$ 是 $P(x, \sigma_k)$ 的全局极小点，则 $\{x_k\}$ 的每个极限点都是原COP问题的全局极小点。
\end{theorem}

\begin{notebox}[外点罚函数法的优缺点]
\textbf{优点}：不要求初始点为可行点

\textbf{缺点}：
\begin{itemize}
    \item $\sigma \to \infty$ 时Hessian条件数恶化，数值不稳定
    \item 中间迭代点不可行
    \item 需要求解一系列越来越困难的无约束问题
\end{itemize}
\end{notebox}

\subsection{内点法（障碍函数法）}

\begin{conceptbox}[内点法思想]
在可行域\textbf{内部}设置\textbf{障碍函数}（Barrier Function），使迭代点始终保持在可行域内部。当点接近边界时，障碍函数趋向 $+\infty$。

\textbf{障碍问题}：
\[
\min_x\ B(x, \mu) = f(x) + \mu \cdot b(x)\quad \text{s.t.}\quad x \in \text{int}(\mathcal{X})
\]
其中 $\mu > 0$ 为障碍参数，随迭代减小至0。
\end{conceptbox}

\textbf{常用障碍函数}：

\textbf{倒数障碍函数}：
\[
b(x) = \sum_{i=1}^{m} \frac{-1}{g_i(x)}
\]

\textbf{对数障碍函数}（更常用）：
\[
b(x) = -\sum_{i=1}^{m} \ln(-g_i(x))
\]

\begin{algobox}[内点法算法]
\textbf{Step 1}：给定\textbf{严格内点} $x_0$（$g_i(x_0) < 0$），$\mu_0 > 0$，缩减因子 $\beta \in (0, 1)$，$\varepsilon > 0$，$k = 0$

\textbf{Step 2}：以 $x_k$ 为初值，求解 $\min_x B(x, \mu_k)$，得 $x_{k+1}$

\textbf{Step 3}：若 $\mu_k \cdot m < \varepsilon$（$m$ 为约束数），停

\textbf{Step 4}：$\mu_{k+1} = \beta \mu_k$，$k \leftarrow k+1$，转Step 2
\end{algobox}

\begin{theorem}[内点法收敛性]
设 $\mu_k \to 0$，$x_k$ 是障碍问题的解，则 $\{x_k\}$ 的每个极限点都是原问题的最优解。
\end{theorem}

\begin{example}[内点法——迭代两次]\label{ex:interior}
$\min\ f(x) = x^2$，s.t. $x \geq 2$。用对数障碍函数，$x_0 = 3$，$\mu_0 = 1$。

\textbf{解}：约束改写为 $g(x) = 2 - x \leq 0$。

障碍函数：$B(x, \mu) = x^2 - \mu \ln(x - 2)$

求导：$\frac{\partial B}{\partial x} = 2x - \frac{\mu}{x-2} = 0 \Rightarrow 2x(x-2) = \mu$

\textbf{第0轮}：$\mu_0 = 1$，解 $2x^2 - 4x - 1 = 0$：
\[
x_1 = \frac{4 + \sqrt{16+8}}{4} = \frac{4 + 2\sqrt{6}}{4} \approx 2.2247
\]

\textbf{第1轮}：$\mu_1 = 0.5$，解 $2x^2 - 4x - 0.5 = 0$：
\[
x_2 = \frac{4 + \sqrt{16+4}}{4} = \frac{4 + 2\sqrt{5}}{4} \approx 2.1180
\]

随着 $\mu \to 0$，$x_k \to 2 = x^*$。
\end{example}

\section{增广拉格朗日函数法与乘子法}\label{sec:augmented-lagrangian}

\subsection{罚函数的困难}

外点罚函数法要求 $\sigma \to \infty$ 才能收敛到可行解，这导致：
\begin{itemize}
    \item Hessian矩阵条件数恶化
    \item 子问题求解越来越困难
    \item 数值不稳定
\end{itemize}

\begin{notebox}[关键观察]
在罚函数最优解 $x^*(\sigma)$ 处：
\[
\nabla P(x^*, \sigma) = \nabla f(x^*) + \sigma \sum_{i} g_i^+(x^*) \nabla g_i(x^*) = 0
\]
对比KKT条件：$\nabla f(x^*) + \sum_i \mu_i^* \nabla g_i(x^*) = 0$

要使两者一致，需要 $\sigma g_i^+(x^*) \approx \mu_i^*$，即约束违反度 $g_i^+(x^*) \approx \mu_i^*/\sigma$。当 $\sigma \to \infty$ 时违反度趋于0，但中间过程约束不被精确满足。
\end{notebox}

\subsection{增广拉格朗日函数法}

\begin{conceptbox}[增广拉格朗日函数]
将罚函数与拉格朗日函数结合，对\textbf{等式约束}问题 $h_j(x) = 0$：
\[
\mathcal{L}_A(x, \lambda, \sigma) = f(x) + \sum_j \lambda_j h_j(x) + \frac{\sigma}{2}\sum_j h_j(x)^2
\]
其中 $\lambda$ 是拉格朗日乘子估计。
\end{conceptbox}

\begin{algobox}[增广拉格朗日法（乘子法）]
\textbf{Step 1}：给定 $x_0$，$\lambda_0$，$\sigma_0 > 0$，$k = 0$

\textbf{Step 2}：求解 $\min_x \mathcal{L}_A(x, \lambda_k, \sigma_k)$ 得 $x_{k+1}$

\textbf{Step 3}：更新乘子：
\[
\lambda_{k+1} = \lambda_k + \sigma_k h(x_{k+1})
\]

\textbf{Step 4}：若 $\|h(x_{k+1})\| < \varepsilon$，停

\textbf{Step 5}：可适当增大 $\sigma_{k+1}$，$k \leftarrow k+1$，转Step 2
\end{algobox}

\begin{notebox}[增广拉格朗日法的优势]
\begin{itemize}[itemsep=0.3em]
    \item 若 $\lambda$ 接近最优乘子 $\lambda^*$，有限 $\sigma$ 即可使约束精确满足
    \item 避免了 $\sigma \to \infty$ 导致的数值困难
    \item 乘子更新使得约束违反度与 $\mu$ 成正比，接近最优时自动减小
    \item 收敛速度比纯罚函数法快得多
\end{itemize}
\end{notebox}

\subsection{不等式约束的处理}

引入松弛变量 $s_i \geq 0$ 将不等式转为等式：$g_i(x) + s_i = 0$。

增广拉格朗日函数变为：
\[
\mathcal{L}_A(x, s, \mu, \sigma) = f(x) + \sum_i \mu_i(g_i(x) + s_i) + \frac{\sigma}{2}\sum_i (g_i(x) + s_i)^2
\]

\begin{example}[增广拉格朗日法——迭代两次]\label{ex:alm}
$\min\ f(x) = x^2$，s.t. $x = 1$。

\textbf{解：}

增广拉格朗日函数：$\mathcal{L}_A = x^2 + \lambda(x-1) + \frac{\sigma}{2}(x-1)^2$

\textbf{参数}：$\lambda_0 = 0$，$\sigma = 2$

\textbf{第0轮}：$\min_x\ x^2 + (x-1)^2 = 2x^2 - 2x + 1$
\[
\frac{d}{dx} = 4x - 2 = 0 \Rightarrow x_1 = 0.5
\]
更新：$\lambda_1 = \lambda_0 + \sigma(x_1 - 1) = 0 + 2(-0.5) = -1$

\textbf{第1轮}：$\mathcal{L}_A = x^2 - (x-1) + (x-1)^2 = x^2 - x + 1 + x^2 - 2x + 1 = 2x^2 - 3x + 2$
\[
\frac{d}{dx} = 4x - 3 = 0 \Rightarrow x_2 = 0.75
\]
更新：$\lambda_2 = -1 + 2(-0.25) = -1.5$

\textbf{第2轮}：$\mathcal{L}_A = x^2 - 1.5(x-1) + (x-1)^2 = 2x^2 - 3.5x + 2.5$
\[
x_3 = \frac{3.5}{4} = 0.875
\]
可见 $x_k \to 1$，$\lambda_k \to -2 = \lambda^*$。
\end{example}

\section{线性规划的内点法}\label{sec:interior-point-lp}

\subsection{解析中心}

\begin{definition}[解析中心（Analytic Center）]
设集合 $S$ 由不等式定义：$S = \{x \in \R^n \mid a_i^T x \geq b_i,\ i = 1, \ldots, m\}$，内点集非空。定义\textbf{势函数}：
\[
\psi(x) = -\sum_{i=1}^{m} \ln(a_i^T x - b_i)
\]
集合 $S$ 的\textbf{解析中心}定义为：
\[
x_{AC} = \argmin_{x \in \text{int}(S)} \psi(x)
\]
\end{definition}

\begin{property}[解析中心性质]
\begin{itemize}[itemsep=0.3em]
    \item 解析中心是唯一的（因为 $\psi$ 严格凸）
    \item 解析中心与不等式的表示方式有关（冗余约束会改变中心）
    \item 势函数是松弛变量对数和的负数，等价于松弛变量乘积的倒数的对数
\end{itemize}
\end{property}

\begin{example}[单位立方体的解析中心]
$S = \{x \in \R^n \mid x_i \geq 0,\ 1 - x_i \geq 0,\ i = 1, \ldots, n\}$

\textbf{解：}

势函数：$\psi(x) = -\sum_{i=1}^{n}[\ln(x_i) + \ln(1-x_i)]$

求导：$\frac{\partial \psi}{\partial x_i} = -\frac{1}{x_i} + \frac{1}{1-x_i} = 0 \Rightarrow x_i = \frac{1}{2}$

解析中心为立方体中心 $x_{AC} = (\frac{1}{2}, \ldots, \frac{1}{2})^T$。
\end{example}

\subsection{中心路径}

考虑LP问题：$\min\ c^T x$ s.t. $Ax = b$，$x \geq 0$。

\begin{definition}[中心路径]
对 $\mu > 0$，定义\textbf{障碍问题}：
\[
\min\ c^T x - \mu \sum_{j=1}^{n} \ln(x_j)\quad \text{s.t.}\quad Ax = b,\ x > 0
\]
当 $\mu$ 从 $\infty$ 变化到0时，障碍问题最优解的轨迹形成\textbf{中心路径}（Central Path），连接解析中心（$\mu \to \infty$）和LP最优面（$\mu \to 0$）。
\end{definition}

\subsection{KKT条件与原始-对偶系统}

引入拉格朗日乘子 $y$（对偶变量），障碍问题的KKT条件：
\begin{align*}
c - \mu X^{-1}e - A^T y &= 0 \\
Ax - b &= 0
\end{align*}

令 $s = \mu X^{-1}e$（即 $Xs = \mu e$），条件变为：
\[
\begin{cases}
Xs = \mu e \\
Ax = b \\
A^T y + s = c
\end{cases}
\]

\begin{definition}[原始-对偶内点法]
求解一系列 $\mu_k \to 0$ 的扰动KKT系统，用Newton法求解每个子问题。
\end{definition}

\subsection{Newton步的求解}

对当前点 $(x, y, s)$（满足 $x > 0, s > 0$），求Newton步 $(\Delta x, \Delta y, \Delta s)$：
\[
\begin{pmatrix}
S & 0 & X \\
A & 0 & 0 \\
0 & A^T & I
\end{pmatrix}
\begin{pmatrix} \Delta x \\ \Delta y \\ \Delta s \end{pmatrix} =
\begin{pmatrix} \mu e - Xs \\ b - Ax \\ c - A^T y - s \end{pmatrix}
\]

\textbf{简化计算}：
\[
A S^{-1} X A^T \Delta y = b - Ax + A S^{-1}(X(c-A^T y-s) - \mu e + Xs)
\]

\begin{notebox}[预测-校正方法（Predictor-Corrector）]
Mehrotra预测-校正是最常用的原始-对偶内点法：
\begin{enumerate}[label=\arabic*.]
    \item \textbf{预测步}（仿射步）：求解 $\mu = 0$ 的Newton系统，预测中心路径方向
    \item \textbf{计算步长}：确定最大可行步长 $\alpha_P, \alpha_D$
    \item \textbf{更新} $\mu$：根据互补性减少程度调整
    \item \textbf{校正步}：加入中心性校正，求解新的Newton系统
\end{enumerate}
\end{notebox}

\section{ADMM与邻近算法简介}\label{sec:admm-proximal}

\subsection{ADMM方法}

\begin{definition}[交替方向乘子法（ADMM）]
考虑结构化优化问题：
\[
\min\ f(x) + g(z)\quad \text{s.t.}\quad Ax + Bz = c
\]

ADMM迭代格式：
\begin{align*}
x_{k+1} &= \argmin_x\ \mathcal{L}_\rho(x, z_k, y_k) \\
z_{k+1} &= \argmin_z\ \mathcal{L}_\rho(x_{k+1}, z, y_k) \\
y_{k+1} &= y_k + \rho(Ax_{k+1} + Bz_{k+1} - c)
\end{align*}
其中增广拉格朗日函数：
\[
\mathcal{L}_\rho(x, z, y) = f(x) + g(z) + y^T(Ax + Bz - c) + \frac{\rho}{2}\|Ax + Bz - c\|^2
\]
\end{definition}

\begin{notebox}[ADMM的特点]
\begin{itemize}[itemsep=0.3em]
    \item 将复杂问题分解为两个（或多个）子问题分别求解
    \item $x$ 更新和 $z$ 更新可以独立进行，适合并行计算
    \item 广泛用于统计学习、图像处理、分布式优化
    \item 在凸问题上有收敛保证
\end{itemize}
\end{notebox}

\begin{example}[LASSO的ADMM求解]
$\min\ \frac{1}{2}\|Ax - b\|^2 + \lambda\|z\|_1$ s.t. $x - z = 0$

\textbf{解：}

\textbf{ADMM迭代}：
\begin{align*}
x_{k+1} &= (A^T A + \rho I)^{-1}(A^T b + \rho(z_k - u_k)) \\
z_{k+1} &= S_{\lambda/\rho}(x_{k+1} + u_k) \quad \text{（软阈值）} \\
u_{k+1} &= u_k + x_{k+1} - z_{k+1}
\end{align*}
其中 $S_\kappa(v) = \text{sign}(v)\max\{|v| - \kappa, 0\}$ 是软阈值算子。
\end{example}

\subsection{邻近算子与邻近算法}

\begin{definition}[邻近算子（Proximity Operator）]
函数 $f$ 的\textbf{邻近算子}定义为：
\[
\text{prox}_f(x) = \argmin_y\ \left\{f(y) + \frac{1}{2}\|y - x\|^2\right\}
\]
\end{definition}

\begin{property}[邻近算子性质]
\begin{enumerate}[label=(\arabic*)]
    \item 若 $f$ 为闭凸函数，则 $\text{prox}_f(x)$ 对任意 $x$ 有唯一解
    \item $\text{prox}_f(x) = x$ 当且仅当 $x \in \argmin f$
    \item Moreau分解：$x = \text{prox}_f(x) + \text{prox}_{f^*}(x)$
\end{enumerate}
\end{property}

\begin{example}[常见邻近算子]
\textbf{解：}

\begin{itemize}[itemsep=0.3em]
    \item $f(x) = I_C(x)$（指示函数）：$\text{prox}_f = \Pi_C$（投影到 $C$）
    \item $f(x) = \lambda\|x\|_1$：$\text{prox}_f(x) = S_\lambda(x)$（软阈值）
    \item $f(x) = \frac{\lambda}{2}\|x\|^2$：$\text{prox}_f(x) = \frac{x}{1+\lambda}$
\end{itemize}
\end{example}

\begin{algobox}[邻近梯度法（Proximal Gradient Method）]
用于求解 $F(x) = f(x) + g(x)$，其中 $f$ 光滑、$g$ 可能非光滑：
\[
x_{k+1} = \text{prox}_{\alpha_k g}\left(x_k - \alpha_k \nabla f(x_k)\right)
\]
即先沿梯度方向走一步，再做邻近投影。
\end{algobox}

\section{本章小结}\label{sec:ch7-summary}

\subsection{本章知识框架}

\begin{conceptbox}[约束问题的最优化方法——核心内容]
\begin{enumerate}[label=\arabic*., leftmargin=2em]
    \item \textbf{非线性最小二乘}
    \begin{itemize}
        \item 高斯-牛顿法：忽略二阶项，利用 $J^T J$
        \item Levenberg-Marquardt法：$J^T J + \lambda I$，最常用
    \end{itemize}
    \item \textbf{最优性条件}
    \begin{itemize}
        \item 无约束：$\nabla f = 0$，$\nabla^2 f \succeq 0$
        \item 等式约束：$\nabla f + \lambda^T \nabla h = 0$
        \item 一般约束：KKT条件（平稳性+可行性+对偶性+互补性）
    \end{itemize}
    \item \textbf{KKT条件的应用}
    \begin{itemize}
        \item 一阶必要条件和二阶充分条件
        \item 凸规划中KKT是充要条件
    \end{itemize}
    \item \textbf{可行方向法}
    \begin{itemize}
        \item 线性约束：求解LP找可行下降方向
        \item 非线性约束：利用梯度近似
    \end{itemize}
    \item \textbf{制约函数法}
    \begin{itemize}
        \item 外点罚函数法：$\sigma \to \infty$，不要求初始可行
        \item 内点法（障碍函数法）：$\mu \to 0$，要求严格内点
    \end{itemize}
    \item \textbf{增广拉格朗日函数法}（见第~\ref{chap:constrained-theory}~章第~\ref{sec:ch6-penalty}~节理论）
    \begin{itemize}
        \item 结合罚函数和拉格朗日函数
        \item 有限 $\sigma$ 即可精确满足约束
        \item 乘子更新保证收敛
    \end{itemize}
    \item \textbf{线性规划的内点法}
    \begin{itemize}
        \item 解析中心、中心路径
        \item 原始-对偶方法、预测-校正
    \end{itemize}
    \item \textbf{ADMM与邻近算法}
    \begin{itemize}
        \item ADMM：分解复杂问题，交替求解
        \item 邻近算子：软阈值、投影等
        \item 邻近梯度法：$x_{k+1} = \text{prox}_{\alpha g}(x_k - \alpha \nabla f)$
    \end{itemize}
\end{enumerate}
\end{conceptbox}

\subsection{方法特性对比}

\begin{table}[htbp]
\centering
\caption{约束优化方法对比}
\begin{tabular}{@{}lllll@{}}
\toprule
\textbf{方法} & \textbf{约束类型} & \textbf{需可行点} & \textbf{收敛性} & \textbf{特点} \\
\midrule
可行方向法 & 线性/非线性 & 是 & 保证 & 需解方向LP \\
外点罚函数 & 等式+不等式 & 否 & $\sigma\to\infty$ & 数值不稳定 \\
内点法 & 不等式 & 严格内点 & $\mu\to 0$ & 路径跟踪 \\
增广拉格朗日 & 等式+不等式 & 否 & 有限$\sigma$ & 最好平衡 \\
ADMM & 可分结构 & 否 & 凸保证 & 分布式 \\
邻近梯度 & 非光滑项 & 否 & 次线性到线性 & 稀疏优化 \\
\bottomrule
\end{tabular}
\end{table}

\subsection{关键公式汇总}

\begin{table}[htbp]
\centering
\caption{核心公式速查}
\begin{tabular}{@{}ll@{}}
\toprule
\textbf{方法/概念} & \textbf{公式} \\
\midrule
高斯-Newton步 & $(J^T J) d = -J^T r$ \\
LM步 & $(J^T J + \lambda I) d = -J^T r$ \\
KKT平稳性 & $\nabla f + \mu^T \nabla g + \lambda^T \nabla h = 0$ \\
互补松弛 & $\mu_i g_i(x) = 0$ \\
外点罚函数 & $P = f + \frac{\sigma}{2}\sum h_j^2 + \frac{\sigma}{2}\sum (g_i^+)^2$ \\
对数障碍函数 & $B = f - \mu\sum \ln(-g_i)$ \\
增广拉格朗日 & $\mathcal{L}_A = f + \lambda^T h + \frac{\sigma}{2}\|h\|^2$ \\
乘子更新 & $\lambda_{k+1} = \lambda_k + \sigma h(x_{k+1})$ \\
中心路径 & $Xs = \mu e$, $Ax = b$, $A^T y + s = c$ \\
ADMM $x$更新 & $x^+ = \argmin_x \mathcal{L}_\rho(x, z, y)$ \\
ADMM $z$更新 & $z^+ = \argmin_z \mathcal{L}_\rho(x^+, z, y)$ \\
邻近算子 & $\text{prox}_f(x) = \argmin_y\{f(y) + \frac{1}{2}\|y-x\|^2\}$ \\
近端梯度 & $x^+ = \text{prox}_{\alpha g}(x - \alpha \nabla f)$ \\
\bottomrule
\end{tabular}
\end{table}

\subsection{关键术语中英对照}

\begin{table}[htbp]
\centering
\begin{tabular}{@{}ll@{}}
\toprule
\textbf{中文} & \textbf{英文} \\
\midrule
非线性最小二乘 & Nonlinear Least Squares \\
Jacobi矩阵 & Jacobian Matrix \\
高斯-牛顿法 & Gauss-Newton Method \\
Levenberg-Marquardt & Levenberg-Marquardt Method \\
起作用约束 & Active Constraint \\
可行方向 & Feasible Direction \\
KKT条件 & Karush-Kuhn-Tucker Conditions \\
互补松弛性 & Complementary Slackness \\
罚函数法 & Penalty Function Method \\
障碍函数法 & Barrier Function Method \\
内点法 & Interior Point Method \\
增广拉格朗日法 & Augmented Lagrangian Method \\
解析中心 & Analytic Center \\
中心路径 & Central Path \\
ADMM & Alternating Direction Method of Multipliers \\
邻近算子 & Proximity/Proximal Operator \\
近端梯度法 & Proximal Gradient Method \\
软阈值 & Soft Thresholding \\
\bottomrule
\end{tabular}
\end{table}

\newpage

% ==================== 第八章 复合优化与智能优化 ====================
